\documentclass[11pt,a4paper]{article}
\usepackage[left=2.2cm,right=2.2cm,top=1.7cm,bottom=1.7cm]{geometry}
\usepackage[T1]{fontenc}
\usepackage[utf8]{inputenc}
\usepackage{lmodern}
\usepackage[ngerman]{babel}
\usepackage[hidelinks]{hyperref}
\usepackage{parskip}
\pagestyle{empty}

\begin{document}

\begin{flushright}
Kostiantyn Pysanyi\\
Leipzig, Deutschland\\
\href{mailto:kostiantyn.pysanyi@gmail.com}{kostiantyn.pysanyi@gmail.com}\\
+49 1525 8447880
\end{flushright}

Provinzial Holding Aktiengesellschaft\\
Recruiting-Team\\
Provinzial-Allee 1\\
48159 Münster

\begin{flushright}
28. August 2026
\end{flushright}

\textbf{Bewerbung als AI Developer (all genders)}

Sehr geehrtes Recruiting-Team,

ich bewerbe mich als AI Developer, weil ich gern mit unstrukturierten Daten arbeite und ML-Lösungen am liebsten so weit entwickle, dass sie auch außerhalb eines Experiments zuverlässig funktionieren. Bei Dokumenten reicht ein gutes Modell allein nicht aus: Datenaufbereitung, Bewertung und technische Umsetzung müssen ebenso stimmen. Genau diese Verbindung interessiert mich an der Position. Mit meiner Erfahrung in Text- und Bildklassifikation sowie im Aufbau produktiver ML-Systeme glaube ich, Ihr Team gut unterstützen zu können.

Bei RAYLYTIC entwickelte ich Modelle für medizinische Bilddaten vom Datensatz bis zum Deployment. Eine 2D--3D-Registrierungspipeline zur Bestimmung der Implantatlage setzte ich nahezu vollständig selbstständig um; sie erhöhte die Akzeptanzrate um 45\,Prozent und verringerte den Fehler um 15\,Prozent. Für reproduzierbare Entwicklung nutzte ich unter anderem versionierte Daten, Experiment-Tracking und automatisierte Auswertungen. Direkte Erfahrung mit unstrukturiertem Text sammelte ich zuvor als Research Engineer bei einem Ticket-Routing-Klassifikator: Ich entwarf die Daten- und Trainingsabläufe und vereinfachte sowohl Datensatzaktualisierungen als auch die Bewertung des Modells. Dadurch kenne ich sowohl die Modellseite als auch die praktische Arbeit, die nötig ist, damit aus einem Prototyp eine verlässliche Anwendung wird.

Deutsch lerne ich seit meinem zwölften Lebensjahr und lebe seit sieben Jahren in Deutschland. Deutsch ist nicht meine Muttersprache, doch mit C1-Niveau und bestandenem TestDaF arbeite und kommuniziere ich sicher auf hohem professionellen Niveau. In einem von mir geleiteten BMBF-Projekt verfasste ich Berichte, vermittelte komplexe Forschungsthemen in Präsentationen und diskutierte sie mit Projektpartnern, darunter Ärztinnen und Ärzte der Charité. Auch im täglichen Austausch mit deutschen Kolleg:innen und Freund:innen ist Deutsch für mich seit Jahren selbstverständlich.

Ich freue mich darauf, Ihnen meine Erfahrungen und Motivation in einem persönlichen Gespräch näher zu erläutern.

Mit freundlichen Grüßen\\[0.8em]
Kostiantyn Pysanyi

\end{document}
